\section{Related work}
\label{sec:related}

Every work cited was checked on 6 October 2026 against its official page.
Where only an abstract or index record could be read, the bibliography says so.

\paragraph{Selective prediction and failure detection under shift.}
Liang, Peng and Sun generalise selective classification to reject label- and
covariate-shifted inputs as well as ambiguous in-distribution ones, using
margin-based scores, on non-clinical image and text
benchmarks~\cite{liang2024selective}.
Evaluation practice for failure detection has itself been examined. Jaeger and
colleagues evaluated confidence-scoring methods across failure sources under one
protocol and found that a simple softmax-response baseline performed best
overall~\cite{jaeger2023call}. Traub and colleagues
argued that AURC over-weights low-coverage regions and proposed the area under
the generalised risk--coverage curve (AUGRC)~\cite{traub2024augrc}. In medical
imaging, Bernhardt and colleagues reported that none of the advanced
confidence-scoring methods they benchmarked consistently beat the softmax
baseline for in-domain misclassification detection~\cite{bernhardt2022failure}.
A 2026 preprint benchmarks confidence scores across eight medical imaging tasks
under realistic shifts~\cite{steinmetz2026failure}. These studies characterise
how well confidence ranks errors. They do not take a selective policy frozen at
one hospital to another while manipulating clinical context, and they do not
ask how report-derived label availability enters the risk being transported.

\paragraph{Rejection for multi-label chest radiographs.}
Aperstein and colleagues add per-pathology entropy-based rejection to a
DenseNet-121 multi-pathology classifier and tune rejection thresholds by
quantile calibration. They evaluate on NIH ChestX-ray14, MIMIC-CXR and PadChest
within dataset, cross-dataset on a source excluded from training, and under
leave-one-dataset-out shift in both directions~\cite{aperstein2025rejection}.
The model uses images only. From the published abstract we could not establish
whether rejection thresholds were carried unchanged between datasets.

\paragraph{Clinical context as an input.}
Indication text improves chest-radiograph classification on MIMIC-CXR when
fused with images through a BERT-based multimodal
transformer~\cite{jacenkow2022indication}. Attention refinement with clinical
history~\cite{yang2025mcxnet}, and fusion of clinical indications and vital
signs with multi-view images~\cite{sloan2025clinically}, also report gains on
single-source benchmark splits. Missing-modality architectures treat absent
inputs as an engineering problem~\cite{wang2025missing,wang2025moehealth}.
None of these evaluates a second institution under a frozen abstention policy.

\paragraph{Cross-institution generalisation and label validity.}
Cohen and colleagues argued that cross-dataset failure of chest-radiograph
models reflects shift in labels more than in images~\cite{cohen2020limits}.
Multi-centre benchmarking continues to identify cross-centre generalisation as
open~\cite{cxrlt2026}. Model rankings for acquisition-parameter robustness can
reverse on external data~\cite{farquhar2026acquisition}. Report-derived labels
are imperfect proxies for image findings. Radiologists labelling reports
disagree substantially with radiologists labelling the corresponding
images~\cite{jain2021visualchexbert}, and an expert audit found frequent
disagreement with the public labels of MIMIC-CXR and
CheXpert~\cite{rafferty2025limitations}. Using the same labeller at two sites,
here CheXbert~\cite{smit2020chexbert}, does not imply equal labelling accuracy
or equal label availability at both. The present results show how much the
second of these, availability, can matter.

\paragraph{Gap addressed.}
We found no study that evaluates a selective policy frozen at one hospital, at
a second hospital, while manipulating pre-diagnostic clinical context under
preregistered rules. We also found none that isolates, with predictions and
thresholds held fixed, how label availability and handling determine the
transported selective risk. The components we use are established: patient-clustered
bootstrap, margin-based confidence, temperature scaling, split conformal
prediction, and AURC and AUGRC. The contribution is their combination into a
protocol and the evidence about what such a protocol can and cannot identify
with report-derived labels.
